\chapter{개발 환경 세팅}
\label{ch:setup}

이 장에서는 본격적인 LLM 구현에 앞서 필요한 개발 환경을 세팅한다. Python, PyTorch, 그리고 이 책의 코드 저장소 구조를 살펴보고, 첫 번째 ``Hello, LLM'' 프로그램을 실행해본다.

\section{필요한 도구}

\begin{center}
\begin{tabularx}{\textwidth}{lXl}
\toprule
\textbf{도구} & \textbf{용도} & \textbf{권장 버전} \\
\midrule
Python & 구현 언어 & 3.10 이상 \\
PyTorch & 딥러닝 프레임워크 & 2.0 이상 \\
NumPy & 수치 계산 & 1.24 이상 \\
pytest & 단위 테스트 & 7.0 이상 \\
git & 버전 관리 & 최신 \\
\bottomrule
\end{tabularx}
\end{center}

GPU는 선택 사항이다. 이 책의 모든 코드는 CPU에서 동작하도록 작성되었다. 다만 대규모 학습을 위해서는 CUDA 지원 GPU가 강력히 권장된다. 2권에서 분산 학습을 다룰 때는 다중 GPU 환경을 가정한다.

\section{PyTorch 설치}

PyTorch는 공식 홈페이지의 설치 가이드를 따른다. CPU 전용 버전과 CUDA 버전 중 자신의 환경에 맞는 것을 선택한다.

\begin{lstlisting}[style=bashstyle]
# CPU 전용 (이 책의 모든 예제는 CPU로 동작)
pip install torch --index-url https://download.pytorch.org/whl/cpu

# CUDA 12.1 지원 (GPU 사용 시)
pip install torch --index-url https://download.pytorch.org/whl/cu121
\end{lstlisting}

설치 확인은 다음과 같이 한다.

\begin{lstlisting}[language=Python]
import torch
print(torch.__version__)         # 2.x.x
print(torch.cuda.is_available()) # GPU 사용 가능 여부
\end{lstlisting}

\section{저장소 구조}

이 책의 코드는 다음과 같은 구조를 가진다. 각 장에서 다루는 모듈이 별도 파일로 분리되어 있어, 필요한 부분만 가져다 쓸 수 있다.

\begin{lstlisting}[style=bashstyle]
make-llm-basic/
  src/makellm/
    __init__.py
    tokenizer/              # 3장: 토크나이저
      base.py
      bpe.py
      unigram.py
    model/                  # 4-6장: 모델 구조
      embedding.py
      attention.py
      transformer_block.py
      gpt.py
    training/               # 7장: 학습
      dataset.py
      loss.py
      optimizers.py
      trainer.py
    inference/              # 8장: 추론
      sampler.py
      generate.py
    utils/                  # 공통 유틸
      config.py
      seed.py
      logging.py
  tests/                    # pytest 단위 테스트
  book/                     # 이 책의 LaTeX 소스
\end{lstlisting}

\section{첫 번째 예제: 텐서 다루기}

본격적인 구현에 앞서 PyTorch 텐서의 기본을 복습하자. 텐서는 NumPy 배열과 거의 동일하게 동작하지만, GPU 연산과 자동 미분(autograd)을 지원한다는 점이 다르다.

\begin{lstlisting}[language=Python]
import torch

# 텐서 생성
x = torch.randn(2, 3)         # 표준정규분포에서 샘플링
print(x.shape)                # torch.Size([2, 3])
print(x.dtype)                # torch.float32

# 연산
y = x @ x.T                   # 행렬곱: [2, 3] @ [3, 2] -> [2, 2]
z = x.sum()                   # 모든 원소 합 (스칼라)

# 자동 미분
w = torch.randn(3, requires_grad=True)
loss = (w * w).sum()          # L2 norm^2
loss.backward()               # 역전파
print(w.grad)                 # d(loss)/d(w) = 2*w
\end{lstlisting}

\begin{infobox}[자동 미분 (Autograd)]
PyTorch의 가장 강력한 기능 중 하나. 텐서에 수행하는 모든 연산을 내부적으로 ``계산 그래프''로 기록하고, \code{loss.backward()} 한 번으로 모든 파라미터에 대한 그래디언트를 자동 계산한다. 우리가 직접 미분 식을 유도할 필요가 없다. 이 책의 모든 모델은 이 autograd에 의존한다.
\end{infobox}

\section{재현성 확보}

딥러닝 실험은 무작위성이 많아 재현성이 중요하다. 시드를 고정하면 같은 코드를 여러 번 실행해도 같은 결과가 나온다.

\begin{lstlisting}[language=Python]
from makellm.utils import set_seed
set_seed(42)  # Python, NumPy, PyTorch 시드를 한 번에 설정
\end{lstlisting}

이 함수는 Python의 \code{random}, NumPy의 \code{numpy.random}, 그리고 PyTorch의 \code{torch.manual\_seed}를 모두 호출한다. CUDA가 존재하면 CUDA 시드도 설정한다.

\section{단위 테스트 실행}

각 장의 코드는 pytest 기반 단위 테스트와 함께 제공된다. 코드를 수정한 후에는 반드시 테스트를 실행하여 회귀가 없는지 확인하자.

\begin{lstlisting}[style=bashstyle]
cd make-llm-basic
pytest tests/ -v
\end{lstlisting}

\begin{lstlisting}[style=bashstyle, basicstyle=\ttfamily\small]
========================= 48 passed in 2.47s =========================
\end{lstlisting}

48개 테스트가 모두 통과하면 환경 세팅이 완료된 것이다.

\section{요약}

\begin{itemize}
  \item PyTorch 2.x와 pytest를 설치했다.
  \item 저장소 구조를 이해했다: \code{src/makellm/}에 모듈별로 코드가 분리되어 있다.
  \item 텐서 연산과 autograd의 기본을 복습했다.
  \item 재현성을 위해 \code{set\_seed()}를 사용한다.
  \item \code{pytest tests/}로 단위 테스트를 실행할 수 있다.
\end{itemize}

다음 장에서는 본격적으로 토크나이저를 구현한다. 텍스트를 숫자로 바꾸는 이 첫 번째 변환이 모든 LLM의 출발점이다.
